\begin{abstract}
Fake online reviews distort consumer decisions and are now prohibited by regulators, yet automated detectors are often evaluated in ways that overstate their accuracy. This study tested whether a fine-tuned pretrained transformer detects deceptive hotel reviews more accurately than a simple linear baseline when both are evaluated on identical partitions that hold out entire hotels, and at what computational cost. Using the Deceptive Opinion Spam corpus of 1,600 English reviews of 20 Chicago hotels, a classifier combining term frequency--inverse document frequency features with logistic regression was compared with fine-tuned RoBERTa-base under a fixed protocol: three hotel-disjoint partitions, a hotel-grouped five-fold cross-validation, hotel-level bootstrap intervals, paired randomisation tests, formatting perturbations and timing measurements. On the primary test set, RoBERTa reached a macro-averaged F1 score of 0.899 against 0.892 for the baseline; the difference was not significant ($p=0.74$), and RoBERTa scored slightly lower on the other two partitions. Across all 20 hotels, its gain was 0.013 (95\% confidence interval $-$0.011 to 0.035). RoBERTa ranked reviews better (average precision 0.986 versus 0.954) but flagged more truthful reviews as deceptive, labelled every review truthful when whitespace formatting was altered, and classified reviews about 65 times more slowly. Under leakage-controlled evaluation, a well-specified linear model therefore matches a fine-tuned transformer on this benchmark at a fraction of the cost, and robustness and threshold calibration, not headline accuracy, should guide the choice of detector.
\end{abstract}

\begin{IEEEkeywords}
Deceptive opinion spam; fake review detection; RoBERTa; logistic regression; data leakage; statistical significance testing
\end{IEEEkeywords}
